
% !TEX root = ../Testa/Principale.tex
% LTeX: language=it

\chapter{Conclusioni}\label{secConclusioni}

    Tutt'i risultati illustrati nel precedente capitolo corroborano la prospettiva nodo-agente in un contesto urbano per descrivere sia la distribuzione della taglia tra città che il loro indice di Pareto.

    L'approssimazione di rango unitario dalla \cref{defMatriceGradi} produce dei risultati corrispondenti quasi perfettamente a quelli ottenuti esattamente sia con \(N\) limitato (Valle d'Aosta) che elevato (Italia), confermando \cite{Nurisso2024}.%; nel caso dell'Italia l'approssimazione è lievemente peggiore di quelli precedenti minando cosí la validità dell'\cref{ossBontaApprossimazione}.

    Tutte le regole d'interazione, vale a dire d'emigrazione (\cref{ossUguaglianzaInterazioneEmigrazione}), riescono a riprodurre fedelmente l'indice di Pareto della coda della distribuzione reale, mentre la distribuzione in sé è stata complessivamente ben riprodotta con quasi ogni regola considerata. In ogni caso la lognormale bimodale segue meglio la coda rispetto alla Pareto confermando \cite{Gualandi2019SDC,Gualandi2019HBLD}; tuttavia entrambi gli adattamenti hanno pregi e difetti: da un punto di vista ingegneristico la Pareto è piú facile da trattare avendo solo un parametro, sebbene sia al tempo stesso peggiore matematicamente; viceversa, la lognormale bimodale è piú corretta a livello matematico ma è composta da ben \(5\) parametri, rendendo piú difficile un eventuale confronto.

    In conclusione, è naturale che il lavoro qui presentato non sia esaustivo e conduca a molteplici sviluppi futuri:

    \begin{enumerate}[
        label=\arabic*.,
        % topsep=0.5em,
        % parsep=0em,
        % itemsep=0.25em,
        % leftmargin=2em,
        % rightmargin=1.5em,
        % \leftmargin + \itemindent = \labelindent + \labelwidth + \labelsep
        %itemindent=!,
        %labelindent=3em,
        %labelwidth=!,
        %labelsep=!,
    ]
        \item sperimentare con regole d'interazione contenenti altri coefficienti della teoria dei grafi come, per esempio, la centralità per intermediazione o il coefficiente d'aggregazione;
        \item ricavare l'equazione di Fokker-Planck della \ref{eqFormaDeboleTipoBoltzmannOmogeneoAssimmetricoUrbano} per studiare analiticamente la densità stazionaria, sebbene la natura non lineare e simmetrica delle regole d'emigrazione qui proposte rendano tale obbiettivo alquanto difficile da raggiungere;
        \item abbandonare l'ipotesi semplificativa della rete statica per descrivere una rete dinamica che evolva parallelamente alla popolazione con un possibile confronto con dati reali [se esistenti];
        \item applicare la presente teoria anche a reti europee o internazionali per confermare la sua validità sia in altri contesti che a scale superiori.
    \end{enumerate}
    